\section{Introduction}
\label{chap:introduction}

% TODO（第1段：研究背景）：从 knowledge graphs 和 KGE-based link
% prediction 的大方向开始，说明 link prediction 的基本任务和实际意义。
% 这里只建立研究背景，不提前展开模型公式、数据集或实验流程。

% TODO（第2段：问题引入）：说明 KGE link prediction 通常使用 Hits@10 等
% aggregate metrics 评价总体准确率，但这些指标不能说明总体表现相近的模型
% 是否会对同一个 query 给出相同的 top-10 decision。

% TODO（第3段：已有研究与研究缺口）：引入 predictive multiplicity，说明
% 已有研究已经观察到 query-level prediction instability，但主要在 global
% model level 汇总。由此提出本文的问题：这种不稳定性是否会随不同的
% relations 和 relation patterns 而变化。

% TODO（第4段：本文工作与贡献）：简要说明本文在 FB15k-237 上分析 RotatE
% 和 TransE 的 repeated runs，并比较 relation-level Hits@10、ambiguity 和
% discrepancy。贡献的证据强度要与第5章一致：mapping type 是主要结果；
% inverse 是探索性副结果；symmetry 是弱结果或负结果。不要把本文写成提出
% 新 KGE model 或 multiplicity mitigation method，也不要在这里堆叠具体数字。

% TODO（第5段：论文结构）：在第2至第4章的最终结构稳定后，简要说明各章
% 的功能。避免逐个小节罗列，只说明 background、related work、analysis
% design、experiments 和 conclusion 之间的关系。
